\chapter{Planejamento para a Criptografia Pós-Quântica (PQC)}

\section{Contextualização e Gestão de Risco}
Esta seção introduz o problema e explica por que a transição deve ser planejada imediatamente, abordando o cenário de risco sob a ótica organizacional.

\subsection{A Ameaça Quântica e o Paradigma ``Harvest Now, Decrypt Later''}
\textbf{Conteúdo:} Explicação sobre como adversários já estão coletando tráfego cifrado atual para quebrar a criptografia assim que computadores quânticos criptograficamente relevantes (CRQC) estiverem disponíveis.\\
\textbf{Documento indicado:} \textit{NIST SP 1800-38A (Executive Summary)}. Fornece o panorama do problema e a urgência da migração sob a perspectiva do NCCoE.

\subsection{Mapeamento do Risco Quântico}
\textbf{Conteúdo:} Como integrar a ameaça quântica nos processos de gestão de risco cibernético da organização.\\
\textbf{Documento indicado:} \textit{NIST CSWP 48 (Cybersecurity White Paper)}. Relaciona a migração PQC com as práticas do \textit{NIST Cybersecurity Framework (CSF)}.

\section{Inventário e Descoberta Criptográfica}
Antes de migrar, uma organização precisa saber o que possui. Esta seção detalha as metodologias para rastrear o uso atual de criptografia.

\subsection{Estratégias de Descoberta em Redes e Sistemas}
\textbf{Conteúdo:} Métodos e ferramentas para varrer infraestruturas, encontrar chaves embutidas em código, protocolos obsoletos e certificados ocultos.\\
\textbf{Documento indicado:} \textit{NIST SP 1800-38B (Discovery for PQC Migration)}. Fornece cenários práticos e arquiteturas de referência para a fase de descoberta.

\subsection{Adoção do CBOM (Cryptographic Bill of Materials)}
\textbf{Conteúdo:} O conceito de catalogar a criptografia de forma padronizada para monitorar vulnerabilidades em cadeias de suprimentos de software.\\
\textbf{Documento indicado:} \textit{NIST SP 1800-38B}. Aborda diretamente como gerar e manter inventários de ativos criptográficos.

\section{Agilidade Criptográfica e Novos Padrões}
Aborda o planejamento arquitetural necessário para facilitar atualizações e apresenta as novas ferramentas de proteção.

\subsection{Conceitos de Agilidade e Desafios de Desempenho}
\textbf{Conteúdo:} O que significa projetar um sistema ágil e quais são os impactos reais na rede (latência, tamanho de chaves, largura de banda) ao trocar os algoritmos.\\
\textbf{Documento indicado:} \textit{NIST SP 1800-38C (Interoperability and Performance)}. Apresenta casos de uso reais e testes de impacto.

\subsection{Os Padrões Pós-Quânticos Finais}
\textbf{Conteúdo:} A especificação técnica e a indicação de uso dos algoritmos selecionados.\\
\textbf{Documentos indicados:}
\begin{itemize}
    \item \textbf{FIPS 203 (ML-KEM):} Para troca de chaves criptográficas.
    \item \textbf{FIPS 204 (ML-DSA) e FIPS 205 (SLH-DSA):} Para assinaturas digitais.
    \item \textbf{NIST IR 8413:} Para critérios técnicos do Round 3.
\end{itemize}

\section{Estratégias de Migração e Hibridização}
Foca na parte prática e operacional da transição.

\subsection{Hibridização (Criptografia Clássica e PQC)}
\textbf{Conteúdo:} A abordagem recomendada de usar algoritmos clássicos em conjunto com PQC durante o período de transição.\\
\textbf{Documento indicado:} \textit{NIST SP 1800-38C}. Discute os protocolos que suportam negociação híbrida.

\subsection{Assinaturas Baseadas em Hash para Adoção Imediata}
\textbf{Conteúdo:} Uso de algoritmos de assinatura que já são considerados seguros contra ataques quânticos, úteis para validar firmware.\\
\textbf{Documento indicado:} \textit{NIST SP 800-208}. Padrões stateful hash-based (LMS e XMSS).

\section{Roadmap e Manutenção da Base Criptográfica}
Conclui o planejamento com prazos normativos e a importância de elementos basilares da criptografia.

\subsection{Cronograma e Descontinuação de Algoritmos Atuais}
\textbf{Conteúdo:} Prazos para que algoritmos de chaves públicas atuais (como RSA-2048) deixem de ser considerados seguros.\\
\textbf{Documento indicado:} \textit{NIST SP 800-131A Rev. 3}. Estipula o ciclo de vida e os prazos limite.

\subsection{O Papel Contínuo da Entropia}
\textbf{Conteúdo:} A garantia de que a segurança fundamental ainda depende da geração de números aleatórios de alta qualidade (RNG).\\
\textbf{Documento indicado:} \textit{NIST SP 800-22}. Conecta os testes de aleatoriedade como pré-requisito indispensável.